Para essa aula, o relatório indica baixo engajamento e alta fragmentação da atenção. A sessão durou 70 minutos, com taxa geral de foco de 59,3\% e tempo médio em foco por aluno de apenas 12,4 minutos, revelando atenção pouco sustentada. A presença foi baixa (54,3\%), e o engajamento por câmera e microfone permaneceu nulo (0,0\%), reforçando a falta de interação. Houve 16 trocas de aba/janela e intensidade de multitarefa de 27,0\%, resultando em fragmentação de foco de 610,0\%, o que comprometeu a continuidade do processamento cognitivo.

A Figura~\ref{fig:poo_quartil} mostra a distribuição temporal do engajamento. O pico ocorreu às 10:43, com queda para 50\% do valor máximo logo nos primeiros 10 minutos, sugerindo início desmotivado e maior atenção apenas na segunda metade da aula.

\begin{figure}[H]
    \centering
    \begin{minipage}
        {0.9\textwidth}
        \centering
        \includegraphics[width=.9\textwidth]{espeon/analysis.poo.quartil.png}
        \caption{Distribuição Temporal - Programação Orientada a Objetos}
        \label{fig:poo_quartil}
    \end{minipage}
    \hfill
\end{figure}

Os \textit{scores} gerais reforçam esse cenário: \textit{score} de engajamento de 37{,}0\%, saúde da atenção de 83{,}5\% e risco de distração de 39{,}3\%. Em conjunto, os indicadores revelam dificuldade em manter foco contínuo e alta propensão à dispersão, apesar de momentos pontuais de atenção razoável.
